% !PRINTED 639
\weberpage{639}
%<<CONT
aussetzungen, von denen wir nachträglich das Resultat wieder befreien
werden. Diese Voraussetzungen sind:

\gesperrt{1. dass durch diese Substitution der $\xi$ für die $x$ keine
der Resolventen $(\varepsilon, \xi)$, in der $\varepsilon$ eine
imaginäre $n^{\text{te}}$ Einheitswurzel ist, verschwindet.}

Nach (13), (16), (17), §.184 bekommen dann die Functionen $f_\nu, F_\nu,
\Phi_\nu$ bestimmte von Null verschiedene Werthe und wir machen die
zweite Voraussetzung,

% [sic] Print runs the list through $f_{n-1}$, although the system defined in
% §.184 (5) and used immediately below runs only from $f_0$ through $f_{n-2}$.
\gesperrt{2. dass durch dieselbe Substitution nicht zwei der Functionen
$f_0, f_1 \ldots f_{n-1}$ einander gleich werden.}

Wir nehmen an, es gehe durch die Substitution der $\xi$ für die $x$

\[
f_0, f_1 \ldots f_{n-2}
\]

in

\[
k_0, k_1 \ldots k_{n-2},
\]

so dass also die $k_0, k_1 \ldots k_{n-2}$ von einander verschieden sind,
und

\[
\Phi_0, \Phi_1 \ldots \Phi_{n-2}
\]

in

\[
K_0, K_1 \ldots K_{n-2}
\]

über, wo aber unter den $K$ auch gleiche vorkommen können.

Wegen der Eigenschaft $\delta)$ der Functionen $f_\nu$ sind die Grössen
$k_0, k_1 \ldots k_{n-2}$ die Wurzeln einer cyklischen Gleichung
$(n-1)^{\text{ten}}$ (also geraden) Grades $\psi(u) = 0$, so dass, wenn
$u$ eine Variable bedeutet,

\weq{1}{\psi(u) = (u - k_0)(u - k_1) \cdots (u - k_{n-2})}

eine Function $(n-1)^{\text{ten}}$ Grades von $u$ mit rationalen
Coëfficienten ist. Die Grössen $k_0, k_1 \ldots k_{n-2}$ können rational
durcheinander ausgedrückt werden in der Form

\weq{2}{k_1 = \Theta(k_0), \quad k_2 = \Theta(k_1), \ldots k_{n-1} = \Theta(k_{n-2}), \quad k_0 = \Theta(k_{n-1})} % [sic] Print refers to $k_{n-1}$ although the sequence is introduced only through $k_{n-2}$; preserved.
(§.\,163).

Ebenso können die Grössen $K_\nu$ nach dem Schlusssatze des §.184 durch
die $k_\nu$ ausgedrückt werden, und zwar in der Form

\weq{3}{K_0 = \Phi(k_0), \quad K_1 = \Phi(k_1), \ldots K_{n-2} = \Phi(k_{n-2}),}

worin $\Phi$ eine rationale Function (in $\mathfrak{K}$) bedeutet.

Danach liefert uns die Formel (18), §.184 folgendes Resultat: Wir setzen
zur Abkürzung

\weq{4}{\tau_\nu = \sqrt[n]{k_\nu},}
%CONT>>
